\documentclass[11pt]{article}

\usepackage[margin=1in]{geometry}
\usepackage{amsmath}
\usepackage{amssymb}
\usepackage{amsthm}
\usepackage{enumitem}
\usepackage{booktabs}
\usepackage{array}
\usepackage{xurl}
\Urlmuskip=0mu
\urlstyle{tt}
\usepackage[colorlinks=true,linkcolor=blue,citecolor=blue,urlcolor=blue]{hyperref}

\theoremstyle{plain}
\newtheorem{theorem}{Theorem}[section]
\newtheorem{lemma}[theorem]{Lemma}
\newtheorem{corollary}[theorem]{Corollary}
\newtheorem{proposition}[theorem]{Proposition}
\theoremstyle{definition}
\newtheorem{definition}[theorem]{Definition}
\newtheorem{hypothesis}[theorem]{Hypothesis}
\theoremstyle{remark}
\newtheorem{remark}[theorem]{Remark}

\newcommand{\R}{\mathbb{R}}
\newcommand{\E}{\mathbb{E}}
\newcommand{\Prob}{\mathbb{P}}
\newcommand{\pr}[2]{\langle #1,\,#2\rangle}
\newcommand{\gap}{\operatorname{gap}}
\newcommand{\num}{\operatorname{num}}
\newcommand{\Var}{\operatorname{Var}}
\newcommand{\Lam}{\Lambda}
\newcommand{\hash}[1]{{\footnotesize\path{#1}}}
\newcommand{\lbl}[1]{\textsc{#1}}
\newcommand{\cert}[1]{\par\smallskip\noindent{\normalfont\textit{Certificate.} #1}}

\title{Liu's Second Hypothesis Holds, and the Union-Closed Constant\\
$c'=1-m^\ast=0.38270908791873502993\ldots$ Is Unconditional:\\
a Machine-Checked Proof}
\author{J.\ Chen\thanks{Repository manuscript package:
\texttt{math/uc/H2\_PAPER/}.  Every theorem carries a status label
(\lbl{proved} / \lbl{machine-verified} / \lbl{computational-evidence} /
\lbl{open}), and every machine-assisted theorem cites the artifact and
SHA-256 that certifies it; Section~\ref{sec:verification} explains how to
replay the whole chain.}}
\date{September 2, 2026}

\begin{document}
\maketitle

\begin{abstract}
Liu (arXiv:2306.08824) showed that the union-closed sets conjecture holds with
constant $c'=1-p^\ast x^\ast\approx0.38271$ under two hypotheses: a
positive-semidefiniteness statement (Section V-A) and the assertion that a
nine-parameter optimisation attains its global minimum at an explicit
two-point law (Section V-B, ``Hypothesis~2'').  We prove that the two-protocol
entropy inequality underlying Hypothesis~2 holds for \emph{every} pair of
Borel probability measures on $[0,1]$ and, more generally, for every
conditionally i.i.d.\ coupling, with the uniform constant $M/m^\ast$ and no
mean constraint.  The inequality Liu's Proposition~3 needs therefore holds
outright, both hypotheses are bypassed, and $c'$ is an unconditional lower
bound for the union-closed constant.  Exact identities in the algebra of
pairings reduce the proof to a sufficient pair of inequalities between
explicit continuous kernels on the unit square, $A\ge0$ and $B\ge0$;
$A\ge0$ (together with the
tight inequality $R=A+B\ge0$, which has a nondegenerate interior zero) is
certified by stratified analytic estimates and a finite Arb branch-and-bound
cover, while $B\ge0$ is equivalent to a parameter-free inequality
$\Phi\ge0$ that follows from an exact identity and a one-variable condition
$\kappa\ge2$ with uniform margin.  Every non-elementary step is a
machine-checked artifact with pinned hashes; a replay driver re-runs the
seven artifacts and asserts byte-identity.  The bound is sharp along a
one-parameter family, and Liu's printed decimal $0.382709087918741$ is
incorrect in its last two digits.  The result concerns the constant only and
says nothing about the conjectured value $1/2$.
\end{abstract}

\section{Introduction and main results}\label{sec:intro}

Let $\mathcal F$ be a union-closed family of subsets of a finite set $E$,
with $\mathcal F\neq\emptyset$ and $\mathcal F\neq\{\emptyset\}$.  Frankl's
conjecture asserts that some element of $E$ belongs to at least half of the
members of $\mathcal F$.  Gilmer~\cite{Gil22} proved the first constant bound
by an entropy argument; the sharp form of his inequality gives
$(3-\sqrt5)/2$~\cite{CL22,AHS22,Saw22}, Sawin's max-entropy coupling improves
it, and the best value obtainable along that route, $c^\ast\approx0.3823455$,
was computed by Yu~\cite{Yu23} and Cambie~\cite{Cam22}.  Liu~\cite{Liu23}
introduced conditionally i.i.d.\ couplings and showed that, with the
specific coupling of his Example~5 (with $f(x)=x(1-x)$) mixed with Gilmer's
independent coupling, the constant improves to
\begin{equation}\label{eq:cprime}
  c' \;=\; 1-m^\ast, \qquad m^\ast=p^\ast x^\ast,
\end{equation}
where $x^\ast$ and $p^\ast$ are the root-defined numbers of his equations
(87)--(88) (Section~\ref{sec:constants}), \emph{provided} two hypotheses
hold (his Theorem~13): the positive semidefiniteness of Section V-A
(``Hypothesis~1''), which drives a Krein--Milman reduction of the admissible
couplings to three shared-mass atoms, and the global-minimiser structure of
Section V-B (``Hypothesis~2''), which asserts that the resulting
nine-parameter optimisation (his (81)--(84)) has optimal value at least $1$.

This paper proves the inequality behind Hypothesis~2 in a form strong enough
to make Hypothesis~1 unnecessary.  Throughout, $\log$ is natural,
\begin{equation}\label{eq:defs}
  h(u)=-u\log u-(1-u)\log(1-u),\quad h(0)=h(1)=0,\qquad
  \pi(s,t)=st\bigl(1+(1-s)(1-t)\bigr),\qquad K=h\circ\pi .
\end{equation}
For a Borel probability measure $\mu$ on $[0,1]$ write $M=\pr{\mu}{x}$ for
its mean and $\pr{\mu}{h}=\int h\,d\mu$.

\begin{definition}[conditionally i.i.d.\ coupling]\label{def:ciid}
Let $(\Omega,P)$ be a probability space and $u\mapsto\nu_u$ a Markov kernel
from $\Omega$ to $[0,1]$.  The measure
$\Gamma=\int\nu_u\otimes\nu_u\,dP(u)$ on $[0,1]^2$ is a
\emph{conditionally i.i.d.\ coupling} of its barycentre
$\mu=\int\nu_u\,dP(u)$ with itself.  Its \emph{numerator} at the parameter
$\beta\in(0,1)$ is
\begin{equation}\label{eq:num}
  \num(\Gamma)=(1-\beta)\pr{\mu\otimes\mu}{h(xy)}
  +\beta\int\pr{\nu_u\otimes\nu_u}{K}\,dP(u).
\end{equation}
\end{definition}

\begin{theorem}[mixture theorem; \lbl{proved}, Theorem~\ref{thm:mixture}]\label{thm:A}
Let $\beta=\beta^\ast$ and $m=m^\ast$ be the constants of
Section~\ref{sec:constants}.  For every Borel probability measure $\mu$ on
$[0,1]$ and every conditionally i.i.d.\ coupling $\Gamma$ of $\mu$,
\[
  \num(\Gamma)\;\ge\;\frac{M}{m^\ast}\,\pr{\mu}{h}.
\]
The inequality is an equality along the whole family
$\mu_w=(1-w)\delta_0+w\delta_{x^\ast}$, $0\le w\le1$, $\Gamma=\mu_w\otimes\mu_w$.
\end{theorem}

For the two-component couplings $\Gamma=(1-q)\,\mu\otimes\mu+q\,\nu\otimes\nu$
this is Liu's two-protocol objective with the mean constraint removed:

\begin{theorem}[Hypothesis~2 for general measures; \lbl{proved}, Theorem~\ref{thm:lift}]\label{thm:B}
For all Borel probability measures $\mu,\nu$ on $[0,1]$ and all
$q\in[0,1]$, with $P=(1-q)\mu+q\nu$ and $M=\pr{P}{x}$,
\[
  \gap(\mu,\nu,q):=(1-\beta^\ast)\pr{P\otimes P}{h(xy)}
  +\beta^\ast\bigl[(1-q)\pr{\mu\otimes\mu}{K}+q\pr{\nu\otimes\nu}{K}\bigr]
  -\pr{P}{h}\;\ge\;\frac{M-m^\ast}{m^\ast}\pr{P}{h},
\]
so $\gap\ge0$ whenever $M\ge m^\ast$.  Liu's nine-parameter optimisation
(81)--(84) at $c=c'$ is the case of finitely supported $\mu,\nu$ with three
shared masses; its optimal value is therefore $\ge1$, i.e.\ Hypothesis~2
holds, with equality at his minimiser $q=0$,
$\mu=(1-p^\ast)\delta_0+p^\ast\delta_{x^\ast}$.
\end{theorem}

\begin{theorem}[the constant is unconditional; \lbl{proved}, Theorem~\ref{thm:uc}]\label{thm:C}
Every union-closed family $\mathcal F\notin\{\emptyset,\{\emptyset\}\}$ of
subsets of a finite set has an element contained in at least
$c'\,|\mathcal F|$ of its members, where
\[
  c'=1-m^\ast\in\bigl[0.382709087918735029930312902098972611626381433
  \pm4.95\times10^{-46}\bigr].
\]
\end{theorem}

The proof of Theorem~\ref{thm:C} re-derives Liu's Proposition~3 from his
definitions (Section~\ref{sec:prop3}); the only inputs are Gilmer's
independent protocol and Liu's Example~5, and Theorem~\ref{thm:A} applied
coordinate by coordinate to the couplings those protocols actually induce.
Liu's Theorem~9 (two-mixture sufficiency), Lemma~11 and Theorem~12 (the
Krein--Milman reduction, which is where Hypothesis~1 enters), and Section
V-A are not used.  Theorem~\ref{thm:A} also supplies the strict constant
$C=(1-c)/m^\ast>1$ that Proposition~3 requires for every $c<c'$, with no
limiting argument.

\paragraph{How the proof is organised.}
Section~\ref{sec:constants} fixes notation and proves the exact structure of
the constants: $x^\ast$ is the unique root in $(0,1)$ of
$z^4-2z^3+3z^2-1$, which is precisely the relation $\pi(x,x)=1-x^2$, so
that $K(x^\ast,x^\ast)=h(x^{\ast2})$ and the kernel
$R=D_{m^\ast}/m^\ast$ vanishes to second order at $(x^\ast,x^\ast)$.
Section~\ref{sec:reduction} proves the reduction: three identities in the
algebra of pairings express $\gap$ through the kernels
\[
  P_2,\;Q_2,\;R=P_2+Q_2,\;\varphi=\sqrt{Q_2(s,s)},\;
  A=P_2+\varphi\otimes\varphi,\; B=Q_2-\varphi\otimes\varphi,
\]
and show that $A\ge0$ and $B\ge0$ pointwise on $[0,1]^2$ imply
Theorem~\ref{thm:B}.  Section~\ref{sec:twovar} proves $R\ge0$ and $A\ge0$ by
stratified analytic estimates plus a finite certified cover.
Section~\ref{sec:phi} proves the parameter-free inequality
$\Phi=K_{st}^2-K_{ss}K_{tt}\ge0$, which is equivalent to $B\ge0$, from an
exact identity and a one-variable condition $\kappa\ge2$.
Section~\ref{sec:lift} composes these into Theorem~\ref{thm:B}, and
Section~\ref{sec:mixture} proves Theorem~\ref{thm:A}.
Section~\ref{sec:prop3} proves Theorem~\ref{thm:C}, Section~\ref{sec:value}
discusses the value of $c'$, and Section~\ref{sec:verification} lists the
artifacts, their hashes, the replay procedure, and the scope of the result.

\paragraph{Labels.}
\lbl{proved} means a complete argument is given here or in the cited
artifact's certificate; \lbl{machine-verified} means an identity or
inequality checked by exact symbolic algebra or by Arb ball arithmetic with
the stated precision, on the stated set; \lbl{computational-evidence} means
a floating-point or sampled observation that carries no weight in any proof;
\lbl{open} means not established.  Nothing labelled
\lbl{computational-evidence} is used in the proofs.

\paragraph{Pairings.}
For finite signed Borel measures $\alpha,\beta$ on $[0,1]$ and bounded Borel
$k$ on $[0,1]^2$ we write $\pr{\alpha\otimes\beta}{k}=\iint k(s,t)\,
d\alpha(s)\,d\beta(t)$ and $\pr{\alpha}{f}=\int f\,d\alpha$.  Three facts are
used throughout: bilinearity; symmetry, $\pr{\alpha\otimes\beta}{k}=
\pr{\beta\otimes\alpha}{k}$ whenever $k(s,t)=k(t,s)$; and the product rule
$\pr{\alpha\otimes\beta}{f(s)g(t)}=\pr{\alpha}{f}\pr{\beta}{g}$.  With
subscripts $0,1$ for $\mu,\nu$ we abbreviate
$E_{ab}=\pr{\cdot}{h(xy)}$, $K_{ab}=\pr{\cdot}{K}$, $H_a=\pr{\cdot}{h}$,
$M_a=\pr{\cdot}{x}$, $F_a=\pr{\cdot}{\varphi}$.

\section{Notation, the constants, and the structure at \texorpdfstring{$(x^\ast,x^\ast)$}{(x*,x*)}}\label{sec:constants}

\begin{lemma}[regularity; \lbl{proved}]\label{lem:L0}
$h$ is continuous on $[0,1]$, $0\le h\le\log2$, $h(u)=h(1-u)$, and
$\pi$ maps $[0,1]^2$ into $[0,1]$.  Consequently every kernel appearing in
this paper ($h(xy)$, $K$, $D_M$, $P_2$, $Q_2$, $R$, $\varphi$, $A$, $B$,
$\Phi$) is continuous on the compact square, and every pairing below is
finite.
\end{lemma}
\begin{proof}
$h''(u)=-1/(u(1-u))<0$, $h(1/2)=\log2$, $u\log u\to0$, and the symmetry is
visible from the definition.  With $a=1-s$, $b=1-t$,
$1-\pi(s,t)=a(1-b)+b(1-a)+ab(a+b-ab)$, a sum of nonnegative terms, and
$\pi\ge0$; so $\pi\in[0,1]$.  Compositions and products of continuous
functions are continuous, and $\varphi=\sqrt{\beta K(s,s)}$ is continuous
because $K(s,s)\ge0$.
\end{proof}

\begin{definition}[the constants; Liu (87)--(90)]\label{def:constants}
Let $f(z)=z^4-2z^3+3z^2-1$ and let $x^\ast$ be its unique root in $(0,1)$.
Put
\[
  p^\ast=\frac{h(x^\ast)}{h(x^{\ast2})},\qquad m^\ast=p^\ast x^\ast,\qquad
  c'=1-m^\ast,\qquad
  \beta^\ast=\frac{\bigl(h'(x^\ast)+h(x^\ast)/x^\ast\bigr)/p^\ast-a'}{b'-a'},
\]
where $a'=2x^\ast h'(x^{\ast2})$, $b'=h'(\pi(x^\ast,x^\ast))\,\pi_d'(x^\ast)$,
$\pi_d(x)=\pi(x,x)=x^2(2-2x+x^2)$ and $\pi_d'(x)=2x(2-3x+2x^2)$.
\end{definition}

\begin{lemma}[\lbl{proved}; enclosures \lbl{machine-verified}]\label{lem:constants}
\leavevmode
\begin{enumerate}[label=(\alph*),nosep]
\item $f$ has exactly one root in $(0,1)$: $f(0)=-1<0<f(1)=1$ and
$f'(z)=2z\bigl(2(z-\tfrac34)^2+\tfrac{15}8\bigr)>0$ on $(0,1]$.
\item $f(x)=0$ is the relation $x^2+\pi(x,x)=1$, i.e.\ Liu's (87).  Hence
$\pi(x^\ast,x^\ast)=1-x^{\ast2}$ and $K(x^\ast,x^\ast)=h(x^{\ast2})$.
\item Liu's (88), $p^2h(x^2)-ph(x)=0$ with $p\neq0$, is $p^\ast=h(x^\ast)/h(x^{\ast2})$;
and $\beta^\ast$ is the unique $\beta$ for which the stationarity conditions
(89)--(90) hold at $(p^\ast,x^\ast)$, equivalently (Lemma~\ref{lem:tight}) for
which $\frac{d}{dx}R_\beta(x,x)\big|_{x^\ast}=0$.
\item Certified enclosures (Arb, 400 bits, propagated from a rational bracket
of $x^\ast$ of width $10^{-70}$ on which $f(\mathrm{lo})<0<f(\mathrm{hi})$
is certified):
\begin{align*}
x^\ast&\in[0.690787593924988014150488204636270562580628177\pm9.69\times10^{-47}],\\
p^\ast&\in[0.893604513905465446476770019038769597178677710\pm3.49\times10^{-46}],\\
m^\ast&\in[0.617290912081264970069687097901027388373618567\pm4.95\times10^{-46}],\\
\beta^\ast&\in[0.100052559862893106664628299476947587796206185\pm2.88\times10^{-46}],\\
c'&\in[0.382709087918735029930312902098972611626381433\pm4.95\times10^{-46}].
\end{align*}
\end{enumerate}
\end{lemma}
\begin{proof}
(a) is the displayed computation.  (b): $\pi(x,x)=x^2(1+(1-x)^2)=2x^2-2x^3+x^4$,
so $x^2+\pi(x,x)-1=f(x)$; then $h(1-x^{\ast2})=h(x^{\ast2})$ by the symmetry
of $h$.  (c): the objective at the two-point law $p\delta_x+(1-p)\delta_0$
with $q=0$ is $G(p,x)=(1-\beta)p^2h(x^2)+\beta p^2h(\pi_d(x))-ph(x)$, and
(89)--(90) ask that $dG=0$ along $d(px)=0$, i.e.\ $dp=-(p/x)dx$.  Writing
out $-(p/x)\partial_pG+\partial_xG=0$, using $h(\pi_d(x^\ast))=h(x^{\ast2})$
and $p^\ast h(x^{\ast2})=h(x^\ast)$, one obtains
$(1-\beta)a'+\beta b'=(h'(x^\ast)+h(x^\ast)/x^\ast)/p^\ast$, whose unique
solution is the displayed $\beta^\ast$.  (d) is the content of the
\texttt{constants} blocks of the artifacts of Sections~\ref{sec:twovar},
\ref{sec:lift} and \ref{sec:mixture}, which agree with one another and with
the independent solver \path{uc/liu9_binding.py} (difference enclosures
containing $0$).
\end{proof}

From now on $\beta=\beta^\ast$ and $m=m^\ast$ unless a formula displays
$\beta,m$ as free parameters.  Define, for $(s,t)\in[0,1]^2$ and $M>0$,
\begin{align}
D_M(s,t)&=M\bigl[(1-\beta)h(st)+\beta K(s,t)\bigr]-\tfrac12\bigl(t\,h(s)+s\,h(t)\bigr),\label{eq:DM}\\
P_2(s,t)&=(1-\beta)h(st)-\frac{t\,h(s)+s\,h(t)}{2m},\qquad
Q_2(s,t)=\beta K(s,t)\ge0,\qquad R=P_2+Q_2=\frac{D_m}{m},\label{eq:R}\\
\varphi(s)&=\sqrt{Q_2(s,s)},\qquad A=P_2+\varphi\otimes\varphi,\qquad
B=Q_2-\varphi\otimes\varphi,\qquad R=A+B,\label{eq:AB}\\
\Phi(s,t)&=K(s,t)^2-K(s,s)K(t,t).\label{eq:Phi}
\end{align}
$P_2$ is sign-indefinite (numerically its minimum on the square is
$\approx-0.0693$, on the diagonal near $s=t\approx0.6807$;
\lbl{computational-evidence}, not used); $Q_2\ge0$; $\Phi$ contains no
parameter.

\begin{lemma}[the tight point; \lbl{proved}, Hessian \lbl{machine-verified}]\label{lem:tight}
$R(x^\ast,x^\ast)=0$ and $\nabla R(x^\ast,x^\ast)=0$; the Hessian of $R$ at
$(x^\ast,x^\ast)$ has certified entries
$R_{ss}=R_{tt}\in[0.75319572205793855886\pm10^{-20}]$,
$R_{st}\in[-0.49637762911066539838\pm10^{-20}]$ (the artifact's radii are
$\sim10^{-51}$; every enclosure quoted in this paper is the artifact's ball
rounded to the printed digits with the radius enlarged to cover the
rounding), hence
$\lambda_{\min}\ge0.25681809294727316048$ and the zero is a nondegenerate
interior minimum.  Likewise $A(x^\ast,x^\ast)=0$ and $\nabla A(x^\ast,x^\ast)=0$,
since $A(s,s)=R(s,s)$ on the diagonal and $\partial_sA(s,s)=\partial_sR(s,s)$
for $0<s<1$; the Hessian of $A$ there is different,
$A_{ss}=A_{tt}\in[0.59340358088849405417\pm10^{-20}]$,
$A_{st}\in[-0.33658548794122089368\pm10^{-20}]$, but has the same
certified $\lambda_{\min}\ge0.25681809294727316048$ (the off-diagonal shift
$\varphi'(x^\ast)^2$ cancels in $A_{ss}+A_{st}$).
Moreover $R(0,t)=R(s,0)=0$ and $R(1,t)=(1-\tfrac1{2m})h(t)$ with
$1-\tfrac1{2m}\in[0.19000913473001832145\pm10^{-20}]$.
\end{lemma}
\begin{proof}
On the diagonal, $r(x):=R(x,x)=(1-\beta)h(x^2)+\beta h(\pi_d(x))-x\,h(x)/m$.
At $x^\ast$, Lemma~\ref{lem:constants}(b) gives
$(1-\beta)h(x^{\ast2})+\beta h(\pi_d(x^\ast))=h(x^{\ast2})$ for \emph{every}
$\beta$, and $x^\ast h(x^\ast)/m^\ast=h(x^{\ast2})$ by the definition of
$m^\ast$; so $r(x^\ast)=0$.  Next
$r'(x)=(1-\beta)\,2xh'(x^2)+\beta h'(\pi_d)\pi_d'-(h(x)+xh'(x))/m$, and at
$x^\ast$, with $m^\ast=p^\ast x^\ast$, the condition $r'(x^\ast)=0$ is exactly
$(1-\beta)a'+\beta b'=(h'(x^\ast)+h(x^\ast)/x^\ast)/p^\ast$, which defines
$\beta^\ast$.  Since $R$ is symmetric, $R_s=R_t$ on the diagonal and
$r'=R_s+R_t=2R_s$, so $\nabla R(x^\ast,x^\ast)=0$.  For $A$:
$A(s,s)=P_2(s,s)+\varphi(s)^2=R(s,s)$, and
$\partial_sA(s,s)=\partial_sP_2(s,s)+\varphi'(s)\varphi(s)
=\partial_sP_2(s,s)+\tfrac12\tfrac{d}{ds}Q_2(s,s)=\partial_sP_2(s,s)+\partial_sQ_2(s,s)$
by symmetry of $Q_2$, valid for $0<s<1$ where $\varphi>0$ is smooth.  The
faces: $h(0)=0$, $\pi(0,t)=0$, $\pi(1,t)=t$.
The Hessian entries and the face constant are Arb enclosures evaluated at
the ball of $x^\ast$ (artifact \path{liu9-h2-twovar.json}).
\end{proof}

\begin{remark}[why the constants are exact objects]\label{rem:exact}
Because $R\ge0$ (Theorem~\ref{thm:twovar}) is \emph{tight} at
$(x^\ast,x^\ast)$ with a positive-definite Hessian, the inequality is a
statement about the exact $(\beta^\ast,m^\ast)$: lowering $m$ makes
$R(x^\ast,x^\ast)$ negative to first order in the perturbation, and moving
$\beta$ at fixed $m$ leaves $R(x^\ast,x^\ast)=0$ (Lemma~\ref{lem:tight}) but
destroys the stationarity, so $R$ becomes negative near $(x^\ast,x^\ast)$
(to second order).  A
first version of the artifact of Section~\ref{sec:lift} declared $100$-digit
rational truncations of the constants to be ``the constants''; at those
rationals Liu's minimiser has a \emph{certified negative} gap,
$[-5.9548013684\times10^{-102}\pm10^{-112}]$ (witness W1b of
\path{liu9-h2-general-lift.json}).  All sign claims in this paper are
therefore evaluated at Arb balls that contain the exact constants, and the
enclosures are how those constants are represented, never a licence to
perturb them.
\end{remark}

\section{The reduction: identities in the algebra of pairings}\label{sec:reduction}

Let $\mu,\nu$ be Borel probability measures on $[0,1]$, $q\in[0,1]$,
$P=(1-q)\mu+q\nu$, $M=\pr{P}{x}=(1-q)M_0+qM_1$, and assume $M>0$ (if $M=0$
then $P=\delta_0$ and $\gap=0$).  Let
\[
  I_{ab}=\pr{\cdot\otimes\cdot}{D_M}\ \ (a,b\in\{0,1\}),\qquad
  c=\beta\pr{(\mu-\nu)\otimes(\mu-\nu)}{K}=\beta(K_{00}+K_{11}-2K_{01}),
\]
and $T=2I_{01}/M+c$.

\begin{theorem}[reduction; \lbl{proved}, \lbl{machine-verified}]\label{thm:reduction}
\leavevmode
\begin{enumerate}[label=(C\arabic*),nosep]
\item $\displaystyle \gap(\mu,\nu,q)=\frac{(1-q)^2I_{00}+q^2I_{11}}{M}+q(1-q)\,T$
for every $q\in[0,1]$, with no constraint on $M$.
\item[(C4)] $\displaystyle \frac{D_M}{M}-\frac{D_m}{m}=\frac{(t\,h(s)+s\,h(t))(M-m)}{2Mm}\ge0$
pointwise whenever $M\ge m>0$.
\item[(C5)] At $M=m$: $T_m:=c+2\pr{\mu\otimes\nu}{R}$.
\item[(C6)] $T_m=\pr{\mu\otimes\mu}{B}+\pr{\nu\otimes\nu}{B}+2\pr{\mu\otimes\nu}{A}+(F_0-F_1)^2$.
\end{enumerate}
\cert{\path{liu9-h2-reduction.json}, sha256
\hash{c63ae94e4fbd2397e19ea9e2a4e2c155f1e02265852d0aacad5c97d70f078e3a}
(Arb, 480 bits: (C1) against Liu's audited Matlab transcription
\path{uc/liu9_objective.py} on the nine-variable class, (C4)--(C6), and the
two nine-variable statements (C2) $T=a^{\mathsf T}\mathsf M_Ta$ and (C3)
$\Psi_M(s,t,s,t)=2D_M(s,t)/M$ used only in Remark~\ref{rem:refuted}; four
mutations caught); independently re-derived in a different free-log basis
with exact zero residuals and 4/4 mutations in \path{liu9-psi-reduction.json},
sha256 \hash{9fb5766f1c3110cce9b231d2f2f1b7c1eec67bc3d6ac7602be50fc72b00c07a3};
(C1), (C6) as sympy identities with $\beta,m$ free symbols in
\path{liu9-h2-general-lift.json} (claims L1, L3).}
\end{theorem}
\begin{proof}
(C1).  Expand both sides in the pairings.
$\pr{P\otimes P}{h(xy)}=(1-q)^2E_{00}+2q(1-q)E_{01}+q^2E_{11}$ and
$\pr{P}{h}=(1-q)H_0+qH_1$.  By the product rule,
$\pr{\mu\otimes\nu}{t\,h(s)+s\,h(t)}=M_1H_0+M_0H_1$, so
$I_{01}=M[(1-\beta)E_{01}+\beta K_{01}]-\tfrac12(M_1H_0+M_0H_1)$,
$I_{00}=M[(1-\beta)E_{00}+\beta K_{00}]-M_0H_0$, and similarly $I_{11}$.
On the right-hand side of (C1) the $E$-terms assemble to
$(1-\beta)\pr{P\otimes P}{h(xy)}$; the $K$-terms are
$\beta[(1-q)^2K_{00}+q^2K_{11}+2q(1-q)K_{01}]+q(1-q)\beta[K_{00}+K_{11}-2K_{01}]
=\beta[(1-q)K_{00}+qK_{11}]$; and the $H$-terms are
$-\frac1M[(1-q)H_0((1-q)M_0+qM_1)+qH_1(qM_1+(1-q)M_0)]=-(1-q)H_0-qH_1$.
This is $\gap$.

(C4) is immediate from \eqref{eq:DM}: the bracketed term cancels and
$\tfrac12(t\,h(s)+s\,h(t))(\tfrac1m-\tfrac1M)\ge0$.

(C5).  $2I_{01}/m=2\pr{\mu\otimes\nu}{D_m/m}=2\pr{\mu\otimes\nu}{R}$.

(C6).  $c+2\pr{\mu\otimes\nu}{R}=\beta(K_{00}+K_{11})-2\beta K_{01}+2\pr{\mu\otimes\nu}{P_2}+2\beta K_{01}
=\beta K_{00}+\beta K_{11}+2\pr{\mu\otimes\nu}{P_2}$.  On the other hand, by
the product rule $\pr{\mu\otimes\mu}{\varphi\otimes\varphi}=F_0^2$ etc., so
$\pr{\mu\otimes\mu}{B}+\pr{\nu\otimes\nu}{B}+2\pr{\mu\otimes\nu}{A}+(F_0-F_1)^2
=\beta K_{00}-F_0^2+\beta K_{11}-F_1^2+2\pr{\mu\otimes\nu}{P_2}+2F_0F_1+(F_0-F_1)^2$,
and the $F$-terms cancel.
\end{proof}

\begin{corollary}[\lbl{proved}]\label{cor:reduction}
If $A\ge0$ and $B\ge0$ pointwise on $[0,1]^2$, then for all $\mu,\nu,q$ with
$M\ge m$,
\[
  \gap(\mu,\nu,q)\;\ge\;(1-q)^2\pr{\mu\otimes\mu}{R}+q^2\pr{\nu\otimes\nu}{R}+q(1-q)T_m\;\ge\;0 .
\]
\end{corollary}
\begin{proof}
By (C1), $\gap=(1-q)^2\pr{\mu\otimes\mu}{D_M/M}+q^2\pr{\nu\otimes\nu}{D_M/M}
+q(1-q)[2\pr{\mu\otimes\nu}{D_M/M}+c]$.  Since $(1-q)^2,q^2,q(1-q)\ge0$ and
the measures are positive, (C4) applied pointwise gives the first inequality
(with (C5) for the last bracket).  $R=A+B\ge0$ handles the first two terms,
and (C6) writes $T_m$ as a sum of pairings of $A$ and $B$ against positive
product measures plus a square.
\end{proof}

\begin{remark}[what is false, and why the splitting is rank-one]\label{rem:refuted}
On the nine-variable class, $T=a^{\mathsf T}\mathsf M_Ta$ with
$(\mathsf M_T)_{ij}=P_2(x_i,y_j)+P_2(x_j,y_i)+Q_2(x_i,x_j)+Q_2(y_i,y_j)$ (claim C2
of the reduction artifact), so what is needed is \emph{copositivity} of
$\mathsf M_T$.  Entrywise nonnegativity is false: at $s=\tau=0$ every $Q_2$
term vanishes and the entry reduces to $P_2(t,\sigma)$, with certified
witness $P_2(\tfrac12,\tfrac12)\in[-0.055370810642371949\pm10^{-18}]$.
Identity (C6) sidesteps copositivity: it is Diananda's decomposition of a
copositive form into a nonnegative part ($A$ on cross pairs, $B$ on
same-component pairs) plus a rank-one positive semidefinite part
$(F_0-F_1)^2$, made explicit by the choice $\varphi(s)=\sqrt{Q_2(s,s)}$.  That
choice saturates $B$ on the whole diagonal ($B(s,s)=0$) and is forced at the
diagonal zeros of $R$, i.e.\ at $s\in\{0,x^\ast,1\}$ (the reduction artifact's
\texttt{phi\_sharpness} block records slack elsewhere but no uniform
rescaling that survives), which is why $B\ge0$ had to be proved through the
parameter-free $\Phi$ (Section~\ref{sec:phi}).
\end{remark}

\section{\texorpdfstring{$R\ge0$ and $A\ge0$}{R>=0 and A>=0} on the unit square}\label{sec:twovar}

\begin{theorem}[\lbl{proved} with \lbl{machine-verified} strata]\label{thm:twovar}
At $(\beta^\ast,m^\ast)$, $R(s,t)\ge0$ and $A(s,t)\ge0$ for all
$(s,t)\in[0,1]^2$.
\cert{\path{liu9-h2-twovar.json}, sha256
\hash{a028ff17ae8492211a2ef9de5721de7f3a0afd8bd45bf7d8d9b77b7c90d2e6eb},
internal digest \hash{4ddbf0f38ca881dde204f086076424617edcdbe6a2826dfa945ef963ebdf70c6},
module \path{uc/liu9_h2_twovar_lemmas.py} sha256
\hash{2afe8e242aa5f506c6d54967a133fe246a7e5cb8bad4961b9860e93b86f67163};
Arb at 400 bits; the artifact also certifies
$\Lambda(s,t)=2P_2(s,t)+\beta h(\pi(s,s))+\beta h(\pi(t,t))\ge0$, which is not
used here.}
\end{theorem}

The proof is stratified, because both functions have exact zeros: on the
faces $\{s=0\}\cup\{t=0\}$, at the corner $(1,1)$, and at the interior point
$(x^\ast,x^\ast)$ (Lemma~\ref{lem:tight}).  A function that touches zero cannot
be certified positive by interval subdivision alone, so every zero is handled
analytically and the certified cover runs only on the compact remainder.
Let $F\in\{R,A\}$ and $c_0=1-\beta-\tfrac1{2m}$, $c_1=1-\tfrac1{2m}$; certified
$c_0\in[0.08995657486712521478\pm10^{-20}]$ and
$c_1\in[0.19000913473001832145\pm10^{-20}]$.

\begin{lemma}[faces; \lbl{proved}]\label{lem:faces}
$F(0,t)=F(s,0)=0$; $R(1,t)=R(t,1)=c_1h(t)\ge0$ and $A(1,t)=A(t,1)=c_0h(t)\ge0$.
\end{lemma}
\begin{proof}
$h(0)=0$, $\pi(0,t)=0$ and $\varphi(0)=0$ give the zero faces.  $\pi(1,t)=t$
and $h(1)=0$ give $R(1,t)=(1-\beta)h(t)+\beta h(t)-h(t)/(2m)$;
$\varphi(1)=\sqrt{\beta h(\pi(1,1))}=\sqrt{\beta h(1)}=0$ gives
$A(1,t)=P_2(1,t)=(1-\beta)h(t)-h(t)/(2m)$.
\end{proof}

\begin{lemma}[near-zero strip; \lbl{proved}]\label{lem:strip}
For $0<u<1$, $u\log(1/u)\le h(u)\le u\log(1/u)+u$.  Consequently
\[
  F(s,t)\ \ge\ P_2(s,t)\ \ge\ st\Bigl[c_0\log\frac1{st}-\frac1m\Bigr]
  \qquad(0<s,t\le1),
\]
and if $\min(s,t)\le\varepsilon=10^{-8}$ then
$F(s,t)\ge st\,[c_0\log(1/\varepsilon)-1/m]\ge0$, the bracket being certified
$\ge0.03707961590679940650$.
\end{lemma}
\begin{proof}
The two-sided bound is $0\le-(1-u)\log(1-u)\le u$.  Both $Q_2=\beta K$ and
$\varphi\otimes\varphi$ are nonnegative, so $F\ge P_2$.  Apply the lower bound
to $h(st)$ and the upper bound to $h(s),h(t)$:
$t\,h(s)+s\,h(t)\le st(\log\tfrac1{st}+2)$, whence
$P_2\ge(1-\beta)st\log\tfrac1{st}-\tfrac{st}{2m}(\log\tfrac1{st}+2)$.
If $\min(s,t)\le\varepsilon$ then $st\le\varepsilon$ and, as $c_0>0$,
$c_0\log\frac1{st}\ge c_0\log\frac1\varepsilon$.
\end{proof}

\begin{lemma}[corner $(1,1)$; \lbl{proved}]\label{lem:corner}
Let $a=1-s$, $b=1-t$, $u=a+b\in(0,2\delta]$.  Then, with $\mathcal
A=a\log\tfrac1a+b\log\tfrac1b\le u(\log\tfrac1u+\log2)$,
\[
  R\ \ge\ u\Bigl[\bigl((1-\beta)(1-\tfrac\delta2)+\beta(1-\delta)-\tfrac1{2m}\bigr)\log\tfrac1u-\tfrac{\log2+1}{2m}\Bigr],
  \qquad
  A\ \ge\ u\Bigl[\bigl((1-\beta)(1-\tfrac\delta2)-\tfrac1{2m}\bigr)\log\tfrac1u-\tfrac{\log2+1}{2m}\Bigr].
\]
With $\delta=10^{-4}$ for $R$ and $\delta=10^{-8}$ for $A$ the brackets are
certified $\ge0.24644229080956936883$ and $\ge0.22327437062436302916$
respectively on $u\in(0,2\delta]$, so $F\ge0$ there; $u=0$ is the exact corner
zero.
\end{lemma}
\begin{proof}
$W:=1-st=u-ab$ and $V:=1-\pi(s,t)=W-ab\,st$ satisfy $ab\le u^2/4\le u\delta/2$,
hence $u(1-\delta/2)\le W\le u$ and $u(1-\delta)\le V\le u$.  By the symmetry
$h(z)=h(1-z)$ and the lower bound of Lemma~\ref{lem:strip},
$h(st)=h(W)\ge W\log\tfrac1W\ge u(1-\tfrac\delta2)\log\tfrac1u$ and
$K(s,t)=h(V)\ge u(1-\delta)\log\tfrac1u$.  The subtracted term is
$t\,h(s)+s\,h(t)=(1-b)h(a)+(1-a)h(b)\le h(a)+h(b)\le\mathcal A+u$ by the upper
bound of Lemma~\ref{lem:strip}, and $\mathcal A=u\log\tfrac1u+u\,H_e(a/u)\le
u(\log\tfrac1u+\log2)$ with $H_e$ the binary entropy.  For $A$ drop
$\varphi\otimes\varphi\ge0$ and use only the $h(st)$ term.
\end{proof}

\begin{lemma}[the interior zero; \lbl{proved}, constants \lbl{machine-verified}]\label{lem:local}
Let $z^\ast=(x^\ast,x^\ast)$, let $r_0>0$, and suppose that on the box
$\{z:|z-z^\ast|_\infty\le r_0\}$ all third partial derivatives of $F$ are
bounded in absolute value by $T_3$, and that $\lambda_{\min}(\nabla^2F(z^\ast))\ge\lambda>0$.
Then with $C_3=2\sqrt2\,T_3$ and $\rho=3\lambda/C_3$,
$F(z^\ast+d)\ge\tfrac\lambda2|d|^2-\tfrac{C_3}6|d|^3\ge0$ for all $d$ with
$|d|_2\le\rho$ and $|d|_\infty\le r_0$.  Certified data: for $R$,
$r_0=0.03$, $\lambda=0.25681809294727316048$, $C_3\le28.567618523131339446$,
$\rho\ge0.026969496187369588069$; for $A$, $r_0=0.025$, the same $\lambda$,
$C_3\le32.259197768814229185$, $\rho\ge0.023883243605848030197$.  The squares
$[x^\ast-w,x^\ast+w]^2$ with $w=9/500$ ($R$) and $w=2/125$ ($A$) satisfy
$\sqrt2\,(w+\eta)\le\rho$ (hence also $w+\eta\le r_0$), where
$\eta\le1.19\times10^{-69}$ is the distance from $x^\ast$ to the rational
centre used by the cover; hence $F\ge0$ on those squares.
\end{lemma}
\begin{proof}
$F(z^\ast)=0$ and $\nabla F(z^\ast)=0$ by Lemma~\ref{lem:tight} (exact algebra;
the Arb enclosures of the value and of the gradient components, of radius
below $5\times10^{-68}$ and $10^{-68}$ respectively, contain $0$ and serve as
consistency checks).  Taylor's theorem with Lagrange remainder along the
segment from $z^\ast$ to $z^\ast+d$ gives
$F(z^\ast+d)=\tfrac12d^{\mathsf T}\nabla^2F(z^\ast)d+\tfrac16\sum_{ijk}F_{ijk}(\xi)d_id_jd_k$
with $|\sum_{ijk}F_{ijk}(\xi)d_id_jd_k|\le T_3(|d_1|+|d_2|)^3\le2\sqrt2\,T_3|d|^3$.
The quadratic form is $\ge\lambda|d|^2$, and
$\tfrac\lambda2|d|^2-\tfrac{C_3}6|d|^3\ge0$ iff $|d|\le3\lambda/C_3$.
$\lambda$ is bounded below by $\min(F_{ss}-|F_{st}|,F_{tt}-|F_{st}|)$ evaluated
on the artifact's unrounded Hessian enclosures of $F$ (radius $\sim10^{-51}$;
Lemma~\ref{lem:tight} prints them rounded, which would only give
$\lambda\ge0.25681809294727316046$); $T_3$ is the upper endpoint of the Arb
enclosure of the third partials over the box; $\varphi$ is smooth near
$z^\ast$ because $K(s,s)>0$ for $0<s<1$.
\end{proof}

\begin{proof}[Proof of Theorem~\ref{thm:twovar}]
Lemma~\ref{lem:strip} covers $\{\min(s,t)\le\varepsilon\}$ (and the zero
faces), Lemma~\ref{lem:corner} the corner square $\{a+b\le2\delta\}\supseteq
[1-\delta,1]^2$, and Lemma~\ref{lem:local} the local square around
$z^\ast$.  For each $F$ the remaining set, $[\varepsilon,1]^2$ minus the local
and corner squares, is the union of $14$ exact rational rectangles (the
$4\times4$ grid generated by the cut coordinates, minus the two removed
cells).  On each rectangle a deterministic Arb branch-and-bound certifies
$F>0$: a cell is accepted when one of three certified lower bounds is
positive --- the analytic product bound of Lemma~\ref{lem:strip}, the lower
endpoint of the direct Arb range of $F$ on the cell (with the range of $h$
taken from endpoint values and $\log2$ exactly when the argument interval
crosses $\tfrac12$), or the centred first-order bound
$F(\mathrm{mid})-\sup|F_s|\,w_s-\sup|F_t|\,w_t$ with derivative enclosures over
the whole cell --- and is bisected along its longer side otherwise.  The run
terminates with every cell accepted: for $R$, $7{,}462$ cells processed,
$3{,}738$ accepted, maximal depth $43$, worst certified lower bound at least
$1.01\times10^{-18}$ (on a cell adjacent to the $\varepsilon$-strip); for $A$,
$16{,}928$ processed, $8{,}471$ accepted, depth $48$.  Three (for $R$) and four
(for $A$) mutations of the certificate --- dropping the local stratum,
dropping the corner stratum, understating $C_3$ by a factor $16$, flipping the
sign of the logarithmic term in Lemma~\ref{lem:strip} --- each make the run
fail, as required.
\end{proof}

\section{\texorpdfstring{$\Phi\ge0$}{Phi>=0} on the unit square, hence \texorpdfstring{$B\ge0$}{B>=0}}\label{sec:phi}

Write $\Lam(u)=h(u)/u=\log\tfrac1u+\mu(u)$ for $0<u\le1$, with
$\mu(u)=-(1-u)\log(1-u)/u$, $\mu(0)=1$, $\mu(1)=0$; $\pi_d(x)=\pi(x,x)$;
$g=\Lam\circ\pi_d$ and $\gamma=\log g$ on $(0,1)$; subscripts
$\pi_{st}=\pi(s,t)$, $\Lam_{ss}=\Lam(\pi_{ss})$, etc.; $a=1-s$, $b=1-t$.
Lemmas~\ref{lem:phi-L1} and \ref{lem:phi-L2} are statements on the open
square $0<s,t<1$, where every $\pi$ argument lies in $(0,1]$ and $\Lam$ is
finite; the closed square is handled by Lemma~\ref{lem:phi-L3}.

\begin{lemma}[elementary facts; \lbl{proved}]\label{lem:phi-L0}
\leavevmode
\begin{enumerate}[label=(\alph*),nosep]
\item $u\,h'(u)-h(u)=\log(1-u)$, hence $\Lam'(u)=\log(1-u)/u^2<0$ on $(0,1)$.
\item $\pi_d'(x)=2x(2-3x+2x^2)>0$ on $(0,1]$ (discriminant $-7$) and
$\pi_d''(x)=4(1-3x+3x^2)>0$ (discriminant $-3$): $\pi_d$ and $\pi_d'$ are increasing.
\item $g$ and $\gamma$ are strictly decreasing on $(0,1)$, and
$\gamma'(x)=\log(1-\pi_d(x))\,\pi_d'(x)/(\pi_d(x)\,h(\pi_d(x)))$.
\item $\mu(u)=1-\sum_{k\ge2}u^{k-1}/(k(k-1))$ on $[0,1)$: $\mu$ is
decreasing and concave, $0\le\mu\le1$, continuous on $[0,1]$.
\item $\pi_{ss}\pi_{tt}-\pi_{st}^2=s^2t^2(s-t)^2$.
\item $M_\mu:=\mu(\pi_{ss})+\mu(\pi_{tt})-2\mu(\pi_{st})\le0$.
\item $\Lam\ge0$ on $(0,1]$.
\end{enumerate}
\end{lemma}
\begin{proof}
(a) $h'(u)=\log\frac{1-u}u$, so $uh'-h=u\log(1-u)-u\log u+u\log u+(1-u)\log(1-u)$.
(b) is a direct computation.  (c) chain rule with (a), (b), and
$\Lam(\pi_d)=h(\pi_d)/\pi_d$.  (d) $-(1-u)\log(1-u)=u-\sum_{k\ge2}u^k/(k(k-1))$
from the series of $-\log(1-u)$; all coefficients after the constant are
negative.  (e) $\pi_{ss}=s^2(1+a^2)$, $\pi_{tt}=t^2(1+b^2)$,
$\pi_{st}=st(1+ab)$, and $(1+a^2)(1+b^2)-(1+ab)^2=(a-b)^2=(s-t)^2$.
(f) by (e), $\pi_{st}\le\sqrt{\pi_{ss}\pi_{tt}}\le\tfrac12(\pi_{ss}+\pi_{tt})$;
$\mu$ decreasing gives $\mu(\pi_{st})\ge\mu(\tfrac12(\pi_{ss}+\pi_{tt}))$ and
concavity gives $\mu(\tfrac12(\pi_{ss}+\pi_{tt}))\ge\tfrac12(\mu(\pi_{ss})+\mu(\pi_{tt}))$.
(g) $h\ge0$.
\end{proof}

\begin{lemma}[exact identity; \lbl{proved}]\label{lem:phi-L1}
For $0<s,t<1$, with $\delta=(s-t)^2/(1+ab)^2$ (so $\pi_{ss}\pi_{tt}=\pi_{st}^2(1+\delta)$) and
$\bar\Lam=\tfrac12(\Lam_{ss}+\Lam_{tt})$,
\[
  \Phi=\pi_{st}^2\Bigl[\tfrac14(\Lam_{ss}-\Lam_{tt})^2
  +\tfrac12\bigl(\log(1+\delta)-M_\mu\bigr)(\Lam_{st}+\bar\Lam)\Bigr]
  -s^2t^2(s-t)^2\Lam_{ss}\Lam_{tt},
\]
and therefore, since the middle term is nonnegative,
\begin{equation}\label{eq:dagger}
  \Phi\ \ge\ s^2t^2\Bigl[\tfrac14(1+ab)^2\bigl(g(s)-g(t)\bigr)^2-(s-t)^2g(s)g(t)\Bigr].
\end{equation}
\end{lemma}
\begin{proof}
$K=\pi\Lam(\pi)$, so $\Phi=\pi_{st}^2(\Lam_{st}^2-\Lam_{ss}\Lam_{tt})
-(\pi_{ss}\pi_{tt}-\pi_{st}^2)\Lam_{ss}\Lam_{tt}$, and
$\Lam_{st}^2-\Lam_{ss}\Lam_{tt}=\tfrac14(\Lam_{ss}-\Lam_{tt})^2+(\Lam_{st}-\bar\Lam)(\Lam_{st}+\bar\Lam)$.
Since $\Lam=\log\frac1\cdot+\mu$,
$2(\Lam_{st}-\bar\Lam)=\log\frac{\pi_{ss}\pi_{tt}}{\pi_{st}^2}-M_\mu=\log(1+\delta)-M_\mu$,
and Lemma~\ref{lem:phi-L0}(e) gives $\pi_{ss}\pi_{tt}-\pi_{st}^2=s^2t^2(s-t)^2$
and $\delta=s^2t^2(s-t)^2/\pi_{st}^2$.  The middle term is $\ge0$ by
$\log(1+\delta)\ge0$, (f) and (g).  Dropping it and writing
$\pi_{st}=st(1+ab)$, $\Lam_{ss}=g(s)$, $\Lam_{tt}=g(t)$ gives
\eqref{eq:dagger}.
\end{proof}

\begin{lemma}[master condition; \lbl{proved}]\label{lem:phi-L2}
For $0<s<t<1$ put $\kappa(s,t)=(1+ab)\,\dfrac{\gamma(s)-\gamma(t)}{t-s}>0$.  Then
\[
  \Phi(s,t)\ \ge\ s^2t^2(s-t)^2g(s)g(t)\Bigl(\frac{\kappa^2}4-1\Bigr);
\]
in particular $\kappa\ge2$ implies $\Phi\ge0$, and $\kappa\ge\tfrac94$ implies
$\Phi\ge\tfrac{17}{64}s^2t^2(s-t)^2g(s)g(t)$.
\end{lemma}
\begin{proof}
With $x=\tfrac12(\gamma(s)-\gamma(t))>0$, $g(s)-g(t)=e^{\gamma(s)}-e^{\gamma(t)}
=2\sqrt{g(s)g(t)}\sinh x$, so the bracket in \eqref{eq:dagger} equals
$(s-t)^2g(s)g(t)\bigl[((1+ab)\sinh x/(t-s))^2-1\bigr]$, and $\sinh x\ge x$
gives $(1+ab)\sinh x/(t-s)\ge\kappa/2$.
\end{proof}

\begin{lemma}[boundary and symmetry; \lbl{proved}]\label{lem:phi-L3}
$\Phi(s,s)=0$, $\Phi(0,t)=\Phi(s,0)=0$, $\Phi(s,1)=\Phi(1,s)=h(s)^2\ge0$, and
$\Phi(s,t)=\Phi(t,s)$.  Hence $\Phi\ge0$ on $[0,1]^2$ follows from
$\kappa\ge2$ on the open triangle $0<s<t<1$.
\end{lemma}
\begin{proof}
$\pi(0,t)=0$, $\pi(s,1)=s$, $\pi(1,1)=1$ and $h(0)=h(1)=0$.
\end{proof}

\begin{lemma}[derivative bounds; \lbl{proved}]\label{lem:phi-L45}
\leavevmode
\begin{enumerate}[label=(N\arabic*),start=0,nosep]
\item For $0<x\le x_0<e^{-(1+\log2)/2}$:
$|\gamma'(x)|\ge\nu_0(x_0):=\dfrac{2-3x_0}{x_0\,(2\log(1/x_0)+1-\log2)}$.
\item For $\tfrac34\le x_1\le x<1$, with $\ell_1=\log\frac1{1-\pi_d(x_1)}$:
$|\gamma'(x)|\ge\nu_1(x_1):=\dfrac{\pi_d'(x_1)}{1-\pi_d(x_1)}\cdot\dfrac{\ell_1}{\ell_1+1}$.
\end{enumerate}
\begin{enumerate}[label=(L5),nosep]
\item $N(x):=-\log(1-\pi_d(x))\,\pi_d'(x)$ is positive and strictly increasing
on $(0,1)$, so for $[l,r]\subset(0,1)$:
$\inf_{[l,r]}|\gamma'|\ge N(l)\big/\bigl(\pi_d(r)\max_{[\pi_d(l),\pi_d(r)]}h\bigr)$,
where the maximum of $h$ over an interval is $\log2$ if the interval
contains $\tfrac12$ and the larger endpoint value otherwise.
\end{enumerate}
\end{lemma}
\begin{proof}
(N0) $-\log(1-\pi)\ge\pi$; $\pi_d'(x)\ge2x(2-3x_0)$ for $x\le x_0$;
$h(\pi)\le\pi(\log\frac1\pi+1)$ (Lemma~\ref{lem:strip}); $\pi_d(x)\le2x^2$ and
$u\mapsto u(\log\frac1u+1)$ is increasing on $(0,1)$, so
$h(\pi_d(x))\le2x^2(2\log\frac1x-\log2+1)$; hence
$|\gamma'(x)|\ge\pi_d'/h(\pi_d)\ge(2-3x_0)/(x(2\log\frac1x+1-\log2))$, and the
denominator is increasing for $x<e^{-(1+\log2)/2}$ (its derivative is
$2\log\frac1x-1-\log2$), so it is at most its value at $x_0$.
(N1) $\pi_d\le1$ and $h(\pi)=h(1-\pi)\le(1-\pi)(\log\frac1{1-\pi}+1)$, so
$|\gamma'|\ge\frac{\ell\,\pi_d'}{(1-\pi_d)(\ell+1)}$ with $\ell=\log\frac1{1-\pi_d}$;
on $[\tfrac34,1]$ the factor $2-3x+2x^2$ increases (vertex at $\tfrac34$), so
$\pi_d'\ge\pi_d'(x_1)$; $\ell/(\ell+1)$ increases in $\ell$, and $\ell$ and
$1/(1-\pi_d)$ increase in $x$.
(L5) $N$ is a product of positive increasing functions
(Lemma~\ref{lem:phi-L0}(b)); $|\gamma'(\xi)|=N(\xi)/(\pi_d(\xi)h(\pi_d(\xi)))$ and
$h$ increases on $[0,\tfrac12]$ and decreases after.
\end{proof}

\begin{theorem}[\lbl{proved} with a \lbl{machine-verified} finite cover]\label{thm:phi}
$\Phi(s,t)\ge0$ on $[0,1]^2$; quantitatively,
$\Phi\ge\tfrac{17}{64}s^2t^2(s-t)^2g(s)g(t)$ for $0<s,t<1$.
\cert{\path{liu9-h2-boundary.json}, sha256
\hash{e4d6d96df3ae435800cef5739fb1ae45b98febf539ab418df116ceb03bd481c9},
internal digest \hash{389804b249cd8b16a47f3c6ab64a70ad4a9afe78097b4a8c6739759804c9eaaa},
module \path{uc/liu9_h2_boundary_layer.py} sha256
\hash{fcb3ed55e8bd9ad6d2f180f5ffa79f192dacc940a5c7506eca2a85aa612905e6}
(Arb, 400 bits).  Independent constructive audit, own cover at
$\theta=23/10$: \path{liu9-h2-phi-audit.json}, sha256
\hash{f2359fa3b3b08bab29965253d777dcf32e4b1aee1e0bacbfb0924b86ac163996},
internal digest \hash{a515102c1807849b606bf4a681bd0a67173445b1a22610d397bb779440222cab},
module \path{uc/liu9_h2_phi_audit.py} sha256
\hash{b1d9ed3d9da9e7c99c026843b683447809fb38d28dba4c49a087729bf66cb50c}.}
\end{theorem}
\begin{proof}
By Lemmas~\ref{lem:phi-L2} and \ref{lem:phi-L3} it suffices to certify
$\kappa\ge\theta:=\tfrac94$ on $0<s<t<1$.  For a rational cell
$C=[s_1,s_2]\times[t_1,t_2]$ with $s_1<t_2$, every $(s,t)\in C$ with
$0<s<t<1$ satisfies, by the mean value theorem and
Lemma~\ref{lem:phi-L0}(c),
\begin{itemize}[nosep]
\item[(D)] $\kappa\ge(1+(1-s_2)(1-t_2))\cdot\inf_{\xi\in[s_1,t_2]\cap(0,1)}|\gamma'(\xi)|$,
the infimum being bounded below by the minimum of $\nu_0(\min(t_2,x_0))$ on
$(0,x_0]$, $\nu_1(\max(s_1,x_1))$ on $[x_1,1)$, and (L5) on the middle piece,
whichever apply (with $x_0=\tfrac18$, $x_1=\tfrac78$);
\item[(Q)] if $s_2<t_1$: $\kappa\ge(1+(1-s_2)(1-t_2))\,(\gamma(s_2)-\gamma(t_1))/(t_2-s_1)$.
\end{itemize}
Starting from the single cell $[0,1]^2$, a cell is accepted when the lower
endpoint of the Arb enclosure of (D) or (Q) exceeds $\theta$, discarded
when $s_1\ge t_2$ (it contains no point with $s<t$), and bisected along its
longer side otherwise.  The run terminates: $30{,}974$ cells processed, $792$
accepted by (D), $14{,}687$ by (Q), $17$ discarded, maximal depth $23$, minimum
certified bound $[2.2500008364172621445\pm10^{-19}]>\theta$.  Five
mutations are caught: raising $\theta$ to $12/5$ (above the true minimum of
$\kappa$; witness $\kappa(\tfrac{27}{50},\tfrac{27}{50}+10^{-6})\in
[2.3941933797780204\pm10^{-16}]<12/5$, so the cover cannot be
tightened past the true minimum); dropping the factor $(1+ab)$ (witness
$|\gamma'(\tfrac9{20})|\in[1.9105187072037912\pm10^{-16}]<2$, so the
factor is load-bearing); the wrong protocol $s+t-st$; $\tfrac14$ replaced by
$\tfrac12$ in Lemma~\ref{lem:phi-L1}; and (N0) applied at $x_0=\tfrac12$
beyond its validity.  The audit artifact re-derives Lemma~\ref{lem:phi-L1}
in its own basis (residuals at most $9.4\times10^{-120}$ in Arb and
$1.3\times10^{-101}$ in $100$-digit mpmath at $22$ points down to $2^{-40}$
from the edges and the diagonal), builds an independent cover at
$\theta=23/10$ from a $16\times16$ grid with its own strip lemmas at
$\tfrac1{16}$ and $\tfrac{15}{16}$ ($88{,}390$ cells, depth $17$, accepted
area $2151/4096$ plus discarded $1945/4096$ equal to $1$ exactly, minimum
certified bound $\ge2.30000075$), reproduces both witnesses, and catches
$6/6$ of its own mutations.
\end{proof}

\begin{remark}
$\kappa$ diverges to $+\infty$ at every boundary point of the square
(Lemma~\ref{lem:phi-L45}), so the singularity of $h'$ at entropy arguments
$0$ and $1$ --- which defeats direct interval subdivision of $\Phi$ near the
edges --- never enters: the endpoint strips are the easy part of the cover.
The numerical minimum of $\kappa$ is $\approx2.394$ at $s=t\approx0.54$
(\lbl{computational-evidence}), a uniform margin of $20\%$ over the
threshold $2$.
\end{remark}

\begin{corollary}[$B\ge0$; \lbl{proved}]\label{cor:B}
For every $\beta>0$, $B=\beta\bigl(K_{st}-\sqrt{K_{ss}K_{tt}}\bigr)\ge0$ on $[0,1]^2$.
\end{corollary}
\begin{proof}
$\Phi=(K_{st}-\sqrt{K_{ss}K_{tt}})(K_{st}+\sqrt{K_{ss}K_{tt}})$ with $K\ge0$.
If the second factor is positive, $B=\beta\Phi/(K_{st}+\sqrt{K_{ss}K_{tt}})\ge0$;
if it vanishes then $K_{st}=0=K_{ss}K_{tt}$ and $B=0$.
\end{proof}

\section{Hypothesis 2 for all Borel probability measures}\label{sec:lift}

\begin{theorem}[\lbl{proved}, conditional only on the cited artifacts]\label{thm:lift}
Theorem~\ref{thm:B} holds: for all Borel probability measures $\mu,\nu$ on
$[0,1]$ and all $q\in[0,1]$, $\gap(\mu,\nu,q)\ge\frac{M-m^\ast}{m^\ast}\pr{P}{h}$,
and in particular $\gap\ge0$ when $M\ge m^\ast$.
\cert{\path{liu9-h2-general-lift.json}, sha256
\hash{eb4874d39904593bddc3f2ae638d1bad6756a19d4ebf6e06a60427b6bb7ac1ce},
internal digest \hash{3ff1e797cd5dbcc4273d9f002b9e540b58829d237b2d344f266e0bea46c02bc8},
module \path{uc/liu9_h2_general_lift.py} sha256
\hash{b265e726d110098a09c00de331c4433454973ba741966b3018b77afcc225b5c3}
(Arb 400 bits and sympy; seven mutations caught; pins the artifacts of
Sections~\ref{sec:reduction}--\ref{sec:phi} by hash and claim status and
checks that they use the same definitions of the constants).}
\end{theorem}
\begin{proof}
The statement with $\gap\ge0$ for $M\ge m^\ast$ is Corollary~\ref{cor:reduction}
with Theorem~\ref{thm:twovar} ($A\ge0$, $R\ge0$) and Corollary~\ref{cor:B}
($B\ge0$).  The quantitative form follows from Theorem~\ref{thm:mixture}
below (the two-component case), or directly: by (C4),
$\gap-\frac{M-m}{m}\pr{P}{h}$ equals the expression obtained from
\eqref{eq:DM} with $M$ replaced by $m$ in the bracket, i.e.\ the right-hand
side of Corollary~\ref{cor:reduction}, which is $\ge0$.
Lemma~\ref{lem:L0} guarantees that all pairings are finite and that
integrating a pointwise inequality between continuous kernels against a
positive product measure preserves it; no density or continuity argument in
the measures is needed, so the fact that $h$ is not Lipschitz at $0,1$ is
irrelevant.
\end{proof}

\begin{remark}[Liu's nine-parameter form; product coupling; witnesses]
Liu's (81)--(84) is $\gap\ge0$ restricted to
$\mu=\sum_{i=1}^3a_i\delta_{x_i}$, $\nu=\sum_{i=1}^3a_i\delta_{y_i}$ (shared
masses) with $M\ge1-c$; Hypothesis~2 at $c=c'$ is the statement that the
minimum of the ratio $\num/\pr{P}{h}$ over this class is $\ge1$, and
Theorem~\ref{thm:lift} gives ratio $\ge M/m^\ast\ge1$ for every $\mu,\nu$ with
$\pr{P}{h}>0$ (when $\pr{P}{h}=0$, i.e.\ $P$ is supported on $\{0,1\}$, the
numerator vanishes too and the cross-multiplied inequality reads $0\ge0$).
Conversely, any finitely supported pair $\mu=\sum_ia_i\delta_{x_i}$,
$\nu=\sum_jb_j\delta_{y_j}$ is the paired configuration with shared masses
$g_{ij}=a_ib_j$ and supports $(x_i,y_j)$, because $\gap$ depends on the
configuration only through its marginals (a polynomial identity in the
masses, checked by sympy at seven sizes in the artifact).  The artifact
records four witnesses at the exact balls: W1a, $\gap=0$ at Liu's minimiser
(equality by the algebra of Lemma~\ref{lem:tight}: $D_{m^\ast}(0,\cdot)=0$
and $D_{m^\ast}(x^\ast,x^\ast)=m^\ast h(x^{\ast2})-x^\ast h(x^\ast)=0$); W1b,
the certified negative gap at the truncated constants of
Remark~\ref{rem:exact}; and W2, W3, certified negative gaps
($\gap\le-0.1249$ at $\mu=\nu=\delta_{1/2}$ and $\gap\le-0.1344$ at a
two-atom/one-atom pair with mean $43/90$) showing that the mean constraint
$M\ge m^\ast$ is load-bearing.  A sample of $100$ certified mean-feasible
random pairs (minimum certified gap $\ge0.0058665$) is recorded as
\lbl{computational-evidence} only.
\end{remark}

\section{The mixture theorem}\label{sec:mixture}

\begin{theorem}[\lbl{proved}, identities \lbl{machine-verified}]\label{thm:mixture}
Let $\Gamma=\int\nu_u\otimes\nu_u\,dP(u)$ be a conditionally i.i.d.\ coupling
of $\mu=\int\nu_u\,dP(u)$, and let $M=\pr{\mu}{x}$.  At $(\beta^\ast,m^\ast)$,
\begin{multline*}
  \num(\Gamma)-\pr{\mu}{h}
  \;=\;\iint\pr{\nu_u\otimes\nu_v}{A}\,dP(u)\,dP(v)
  +\int\pr{\nu_u\otimes\nu_u}{B}\,dP(u)\\
  +\Var_P\bigl(\pr{\nu_U}{\varphi}\bigr)
  +\frac{M-m^\ast}{m^\ast}\pr{\mu}{h},
\end{multline*}
and consequently $\num(\Gamma)\ge\frac{M}{m^\ast}\pr{\mu}{h}$.
\cert{\path{liu9-h2-mixture-theorem.json}, sha256
\hash{329f7e2d71af8cd78d1a921c72b9ae4d05b71113134eb3341d69932f19ea24b3},
internal digest \hash{b108b781224dff601ffcf3f562703d3b620e7da6d89cf6d000b6ccbd06a4773b},
module \path{uc/liu9_h2_mixture_theorem.py} sha256
\hash{3098a1ca30a0f16582df02e1fb7dfd8fa27970c031167125e696e58ac43dfbab}:
the identities (I)--(III) below as exact rational-function identities in the
free pairing algebra for $2,3,4$ components with symbolic weights (sympy
residual $0$), in Arb to $10^{-100}$ at $60$ random mixtures with up to $6$
components and up to $4$ atoms each (atoms at $0$ and $1$ included), agreement
with Liu's audited transcription at two components to $10^{-100}$, six
mutations caught.}
\end{theorem}
\begin{proof}
Write $K_{uv}=\pr{\nu_u\otimes\nu_v}{K}$, $P_{uv}=\pr{\nu_u\otimes\nu_v}{P_2}$,
$F_u=\pr{\nu_u}{\varphi}$, $H=\pr{\mu}{h}$.  All integrands below are bounded
and measurable ($u\mapsto\nu_u$ is a Markov kernel and every kernel is
bounded continuous, Lemma~\ref{lem:L0}), so Fubini applies, and by
bilinearity $\pr{\mu\otimes\mu}{k}=\iint\pr{\nu_u\otimes\nu_v}{k}\,dP\,dP$ for
every kernel $k$.  Three identities:
\begin{align*}
\text{(I)}\quad&\int K_{uu}\,dP-\iint K_{uv}\,dP\,dP=\tfrac12\iint(K_{uu}+K_{vv}-2K_{uv})\,dP\,dP,\\
\text{(II)}\quad&\pr{\mu\otimes\mu}{(1-\beta)h(xy)+\beta K}-H=\pr{\mu\otimes\mu}{R}+\tfrac{M-m}{m}H,\\
\text{(III)}\quad&\pr{\mu\otimes\mu}{R}+\tfrac\beta2\iint(K_{uu}+K_{vv}-2K_{uv})\,dP\,dP\\
&\qquad=\iint\pr{\nu_u\otimes\nu_v}{A}\,dP\,dP+\int\pr{\nu_u\otimes\nu_u}{B}\,dP+\Var_P(F_U).
\end{align*}
(I) uses only $\int dP=1$.  (II): $R=(1-\beta)h(xy)+\beta K-\frac{y\,h(x)+x\,h(y)}{2m}$
and, by the product rule, $\pr{\mu\otimes\mu}{y\,h(x)+x\,h(y)}=2MH$.
(III): the left side is $\iint(P_{uv}+\beta K_{uv})\,dP\,dP+\beta\int K_{uu}\,dP
-\beta\iint K_{uv}\,dP\,dP=\iint P_{uv}\,dP\,dP+\beta\int K_{uu}\,dP$; the right
side is $\iint P_{uv}\,dP\,dP+(\int F_u\,dP)^2+\beta\int K_{uu}\,dP-\int F_u^2\,dP
+\int F_u^2\,dP-(\int F_u\,dP)^2$, by the product rule for
$\varphi\otimes\varphi$.  Now
$\num(\Gamma)-H=\pr{\mu\otimes\mu}{(1-\beta)h(xy)+\beta K}-H+\beta\bigl[\int K_{uu}dP-\iint K_{uv}dPdP\bigr]$,
and (I), (II), (III) in turn give the displayed identity.  The first two
terms are $\ge0$ by Theorem~\ref{thm:twovar} and Corollary~\ref{cor:B}
integrated against positive product measures, and the third is a variance.
$R\ge0$ is not used.
\end{proof}

\begin{remark}[two components; the paired chain]
For $P=(1-q)\delta_0+q\delta_1$, $\nu_0=\mu$, $\nu_1=\nu$, the variance is
$q(1-q)(F_0-F_1)^2$ and the channel $\frac\beta2\iint(\cdots)$ is
$q(1-q)\,c$; the displayed identity is then (C1)+(C4)+(C6) of
Theorem~\ref{thm:reduction}, and Theorem~\ref{thm:lift} is the special case.
\end{remark}

\begin{lemma}[sharpness; \lbl{proved}]\label{lem:sharp}
For $\mu_w=(1-w)\delta_0+w\delta_{x^\ast}$, $0\le w\le1$, and the one-component
coupling $\Gamma=\mu_w\otimes\mu_w$: $\num(\Gamma)=\frac{M}{m^\ast}\pr{\mu_w}{h}$
exactly.  Laws supported on $\{0,1\}$ give $\num=\pr{\mu}{h}=0$.  Hence the
constant $M/m^\ast$ in Theorem~\ref{thm:mixture} is attained at every mean
$M\in[0,x^\ast]$, and $c'=1-m^\ast$ is the largest constant obtainable from
this two-protocol inequality at $\beta^\ast$.
\end{lemma}
\begin{proof}
Cross terms vanish because $h(0\cdot x)=0$ and $K(0,x)=h(0)=0$, so
$\num=w^2\bigl[(1-\beta)h(x^{\ast2})+\beta K(x^\ast,x^\ast)\bigr]=w^2h(x^{\ast2})$
by Lemma~\ref{lem:constants}(b).  Also $\pr{\mu_w}{h}=w\,h(x^\ast)$ and
$M=wx^\ast$, so $\frac{M}{m^\ast}\pr{\mu_w}{h}=w^2\frac{x^\ast h(x^\ast)}{m^\ast}=w^2h(x^{\ast2})$.
For laws on $\{0,1\}$ every entropy vanishes.  Arb enclosures of the
residual at the exact balls contain $0$ at $w=\tfrac3{10},\tfrac12,\tfrac9{10},1$.
\end{proof}

\begin{remark}
Liu's minimiser is $w=p^\ast$ ($M=m^\ast$); Lemma~\ref{lem:sharp} shows the
equality is not a property of that single point but of the ray through
$\delta_{x^\ast}$.  The artifact also certifies a five-component mixture with
$M>m^\ast$ whose ratio exceeds $M/m^\ast$ strictly, and, by
Remark~\ref{rem:exact}, that at the $100$-digit truncations the bound is
false at $w=p^\ast$ (certified $-5.95\times10^{-102}$).
\end{remark}

\section{From union-closed families to the constant}\label{sec:prop3}

This section re-derives Liu's Proposition~3 from his Definition~1 and
Example~5 so that the dependence on his paper is limited to definitions.
Every step is elementary; it is labelled \lbl{proved} (human argument,
re-derived independently by two readers, cf.\ Section~\ref{sec:verification}).
Entropies may be taken in any base: the accounting below is base-free, and
the inequality of Theorem~\ref{thm:mixture} is homogeneous in the base.

\paragraph{Setting.}
Identify subsets of $E=\{1,\dots,n\}$ with $0/1$ vectors and let
$X^n=(X_1,\dots,X_n)$ be uniform on $\mathcal F$, so $H(X^n)=\log|\mathcal F|$.
For $i\ge1$ and a past $x^{i-1}$ of positive probability write
$\sigma_i(x^{i-1})=\Prob(X_i=1\mid X^{i-1}=x^{i-1})$.

\begin{definition}[protocol; Liu, Definition~1]
A protocol $\Pi$ assigns to each $(s,t)\in[0,1]^2$ a law $\Pi_{s,t}$ on
$\{0,1\}^2$ with marginal means $s$ and $t$.  Given $\Pi$, generate
$(\tilde X^n,\tilde Y^n)$ sequentially: at step $i$, with
$S_i=\sigma_i(\tilde X^{i-1})$ and $T_i=\sigma_i(\tilde Y^{i-1})$, draw
$(\tilde X_i,\tilde Y_i)\sim\Pi_{S_i,T_i}$ using randomness independent of
everything before.
\end{definition}

\begin{lemma}[\lbl{proved}]\label{lem:protocol}
For every protocol: (i) $\tilde X^n\sim X^n$ and $\tilde Y^n\sim X^n$, and only
pasts of positive probability are visited; (ii) the law of $S_i$ is the law
$\mu_i$ of $\sigma_i(X^{i-1})$, with $\E S_i=\Prob(X_i=1)$; (iii)
\[
  H(\tilde X^n\vee\tilde Y^n)\ \ge\ \sum_{i=1}^n\E\bigl[h\bigl(\Pi_{S_i,T_i}(0,0)\bigr)\bigr],
  \qquad H(X^n)=\sum_{i=1}^n\E\bigl[h(S_i)\bigr],
\]
where $\vee$ is the coordinatewise maximum; and (iv) if $\mathcal F$ is
union-closed then $H(\tilde X^n\vee\tilde Y^n)\le H(X^n)$.
\end{lemma}
\begin{proof}
(i) $\Prob(\tilde X_i=1\mid\tilde X^{i-1},\tilde Y^{i-1})=S_i=\sigma_i(\tilde X^{i-1})$
by the marginal condition, so by induction on $i$ the law of $\tilde X^i$ is
that of $X^i$; likewise for $\tilde Y$.  (ii) follows from (i), and
$\E S_i=\E\,\Prob(X_i=1\mid X^{i-1})=\Prob(X_i=1)$.  (iii) By the chain rule and
because $(\tilde X\vee\tilde Y)^{i-1}$ is a function of
$(\tilde X^{i-1},\tilde Y^{i-1})$,
$H(\tilde X^n\vee\tilde Y^n)=\sum_iH\bigl((\tilde X\vee\tilde Y)_i\mid(\tilde X\vee\tilde Y)^{i-1}\bigr)
\ge\sum_iH\bigl((\tilde X\vee\tilde Y)_i\mid\tilde X^{i-1},\tilde Y^{i-1}\bigr)$;
conditionally on the pasts, $(\tilde X_i,\tilde Y_i)\sim\Pi_{S_i,T_i}$, so
$(\tilde X\vee\tilde Y)_i$ is Bernoulli with
$\Prob(=0\mid\text{pasts})=\Pi_{S_i,T_i}(0,0)$ and the conditional entropy is
$\E[h(\Pi_{S_i,T_i}(0,0))]$, using $h(u)=h(1-u)$.  The second formula is the
chain rule for $X^n$ with (ii).  (iv) $\tilde X^n\vee\tilde Y^n$ is the
indicator of the union of two members of $\mathcal F$, hence takes values in
$\mathcal F$, and the uniform law maximises entropy on a given support.
\end{proof}

\paragraph{The two protocols.}
\begin{enumerate}[label=(\roman*),nosep]
\item Gilmer's protocol $\Pi^{(1)}_{s,t}=\mathrm{Ber}(s)\otimes\mathrm{Ber}(t)$,
$\Pi^{(1)}_{s,t}(0,0)=(1-s)(1-t)$.  The two coordinates use independent
randomness at every step, so $\tilde X^{i-1}$ and $\tilde Y^{i-1}$ are
independent and $(S_i,T_i)\sim\mu_i\otimes\mu_i$.
\item Liu's Example~5 with $f(x)=x(1-x)$.  Let $U$ be uniform on $[0,1]$ and
$\sigma(U)=+1$ if $U>\tfrac12$, $-1$ otherwise.  Given $(s,t)$ and $U$, set
$\tilde X=0$ with probability $(1-s)+f(1-s)\sigma(U)=(1-s)(1+s\sigma(U))$
using $\tilde X$'s local randomness, and $\tilde Y=0$ with probability
$(1-t)(1+t\sigma(U))$ using $\tilde Y$'s local randomness and the same $U$.
These are probabilities ($0\le(1-s)(1\pm s)\le1$), the means are $s$ and $t$
($\E\sigma(U)=0$), and
\[
  \Pi^{(2)}_{s,t}(0,0)=\E\bigl[(1-s)(1+s\sigma)(1-t)(1+t\sigma)\bigr]
  =(1-s)(1-t)(1+st)=\pi(1-s,1-t),
\]
since $\E\sigma=0$ and $\sigma^2=1$.  In the sequential construction a fresh
$U_i$ is drawn at step $i$ and shared by the two coordinates.  Each
$\tilde X_j$ is the same measurable function $G_j$ of
$(\tilde X^{j-1},U_j,\xi_j)$ as $\tilde Y_j$ is of $(\tilde Y^{j-1},U_j,\eta_j)$,
where $(\xi_j)$, $(\eta_j)$ are the local randomness and $(\xi_j)$, $(\eta_j)$,
$(U_j)$ are independent.  Hence, conditionally on $U^{i-1}$, $\tilde X^{i-1}$
and $\tilde Y^{i-1}$ are i.i.d., and so are $S_i=\sigma_i(\tilde X^{i-1})$ and
$T_i=\sigma_i(\tilde Y^{i-1})$: the law of $(S_i,T_i)$ is
$\Gamma_i=\int\nu_u\otimes\nu_u\,dP_{U^{i-1}}(u)$, with $\nu_u$ the regular
conditional law of $S_i$ given $U^{i-1}=u$ and barycentre $\mu_i$.  This is
Liu's observation after his Example~5.
\end{enumerate}

\begin{theorem}[\lbl{proved}]\label{thm:uc}
Theorem~\ref{thm:C} holds.
\end{theorem}
\begin{proof}
If $|\mathcal F|=1$ then $\mathcal F=\{A\}$ with $A\neq\emptyset$ and every
element of $A$ has frequency $1$.  Let $|\mathcal F|\ge2$, so
$H(X^n)>0$.  Fix $c<c'$ and suppose, for contradiction, that
$\Prob(X_i=1)<c$ for every $i$.  Let $\bar\mu_i$, $\bar\Gamma_i$, $\bar\nu_u$
denote the images of $\mu_i$, $\Gamma_i$, $\nu_u$ under $s\mapsto1-s$; then
$\bar\Gamma_i=\int\bar\nu_u\otimes\bar\nu_u\,dP_{U^{i-1}}$ is a conditionally
i.i.d.\ coupling of $\bar\mu_i$, whose mean is
$M_i=\E[1-S_i]=1-\Prob(X_i=1)>1-c$, and $\pr{\bar\mu_i}{h}=\E[h(S_i)]$ by
the symmetry of $h$.  By Lemma~\ref{lem:protocol}(iii) for the two protocols,
\begin{align*}
&(1-\beta^\ast)\,H(\tilde X^{(1)n}\vee\tilde Y^{(1)n})+\beta^\ast H(\tilde X^{(2)n}\vee\tilde Y^{(2)n})\\
&\qquad\ge\sum_i\Bigl[(1-\beta^\ast)\pr{\bar\mu_i\otimes\bar\mu_i}{h(xy)}
+\beta^\ast\!\int\!\pr{\bar\nu_u\otimes\bar\nu_u}{K}\,dP_{U^{i-1}}\Bigr]
=\sum_i\num(\bar\Gamma_i)\\
&\qquad\ge\sum_i\frac{M_i}{m^\ast}\,\E[h(S_i)]
\ \ge\ \frac{1-c}{m^\ast}\sum_i\E[h(S_i)]=\frac{1-c}{m^\ast}\,H(X^n),
\end{align*}
by Theorem~\ref{thm:mixture} for each $i$ and $\E[h(S_i)]\ge0$.  By
Lemma~\ref{lem:protocol}(iv) the left side is at most
$(1-\beta^\ast)H(X^n)+\beta^\ast H(X^n)=H(X^n)$.  Since $H(X^n)>0$ this forces
$(1-c)/m^\ast\le1$, i.e.\ $c\ge1-m^\ast=c'$, contradicting $c<c'$.  Hence for
every $c<c'$ some $i$ has $\Prob(X_i=1)\ge c$; as the $n$ frequencies are
fixed numbers, $\max_i\Prob(X_i=1)\ge c'$, i.e.\ some $i\in E$ lies in at
least $c'|\mathcal F|$ members.
\end{proof}

\begin{remark}[on Liu's argument, and what is not used]
Liu's Proposition~3 asks, for every $c<c'$, for a constant $C>1$ in
$\sum_kw_k\inf_{P_{ST}\in\mathcal C_k(\mu)}\E[h(\Pi^{(k)}_{S,T}(0,0))]\ge C\,\E[h(S)]$
over all laws $\mu$ with mean $\le c$; Theorem~\ref{thm:mixture} gives it with
the explicit $C=(1-c)/m^\ast$, which exceeds $1$ precisely when $c<c'$, so no
limiting or separate strictness argument is needed.  The
couplings that actually occur are the mixtures $\bar\Gamma_i$, so
Theorem~\ref{thm:mixture} is applied to them directly; Liu's
Theorem~9 (reduction to two-point mixing laws), Lemma~11 (the
positive-semidefiniteness of Section V-A, ``Hypothesis~1'') and Theorem~12 (the
Krein--Milman reduction to three shared-mass atoms) are not invoked, and
neither is the weak closure $\mathcal C_3(\mu)$.  Section~V-B's assertion that
the nine-parameter optimum is $\ge1$ is replaced by Theorem~\ref{thm:lift}.
Laws supported on $\{0,1\}$ (his Lemma~7) need no separate treatment: they
contribute $0\ge0$.
\end{remark}

\section{The value of \texorpdfstring{$c'$}{c'}}\label{sec:value}

\begin{theorem}[\lbl{machine-verified}]\label{thm:value}
\[
  c'=1-m^\ast\in\bigl[0.382709087918735029930312902098972611626381433\pm4.95\times10^{-46}\bigr];
\]
in particular $c'>0.38270908791873$.
\cert{\path{liu9-h2-mixture-theorem.json} and \path{liu9-h2-general-lift.json}
(Table~\ref{tab:artifacts}), \texttt{constants.exact\_balls\_45\_digits}.}
\end{theorem}

The enclosure is propagated in Arb from the certified rational bracket of
$x^\ast$ (Lemma~\ref{lem:constants}) through the closed forms of
Definition~\ref{def:constants}, in the artifacts of Sections~\ref{sec:lift}
and \ref{sec:mixture}, and agrees with the independent $100$-digit solver
\path{uc/liu9_binding.py} (difference enclosures containing $0$).  It exceeds
$c^\ast\approx0.3823455$ of \cite{Yu23,Cam22} by $3.6\times10^{-4}$.

Liu's printed values (91)--(94) read $p^\ast\approx0.893604513905457$,
$x^\ast\approx0.690787593924988$, $c'\approx0.382709087918741$,
$\beta^\ast\approx0.100052559862974$.  Of these, (92) agrees with the
certified enclosure to all printed digits, while for (91), (93) and (94)
the printed-minus-certified differences are approximately
$-8.4\times10^{-15}$, $+6.0\times10^{-15}$ and $+8.1\times10^{-14}$
respectively, i.e.\ they are wrong in their last two digits (last three for
$\beta^\ast$); the correct $15$-digit roundings are
$p^\ast\approx0.893604513905465$, $c'\approx0.382709087918735$,
$\beta^\ast\approx0.100052559862893$.  The constants are defined by the
same equations (87)--(90), so the discrepancy is in the printed decimals
only and does not affect any statement above.

\section{Verification: artifacts, replay, labels, scope}\label{sec:verification}

\paragraph{Artifacts.}
Table~\ref{tab:artifacts} lists the seven machine-checked artifacts, all
under \path{math/uc/verification/results/}, produced by the modules under
\path{math/uc/}.  Each artifact is a JSON document with a
\texttt{claim\_status}, an internal digest \texttt{report\_sha256} (the SHA-256
of the canonical compact serialisation of the document with that field
omitted or blanked, in the module's declared scope), and, where the module
records it, the SHA-256 of the module source that produced it.  The
artifacts of Sections~\ref{sec:lift} and \ref{sec:mixture} pin the file
hashes and claim statuses of the artifacts they depend on and fail if they
change; the audit artifact pins the certificate it audits.

\begin{table}[p]
\centering\scriptsize
\begin{tabular}{@{}ll@{}}
\toprule
artifact (module) & claim; status\\
\midrule
\path{liu9-h2-reduction.json} (\path{liu9_h2_reduction.py}) & (C1)--(C6), Theorem~\ref{thm:reduction}; \texttt{REDUCTION-CERTIFIED}\\
\multicolumn{2}{@{}l@{}}{\quad file \texttt{c63ae94e4fbd2397e19ea9e2a4e2c155f1e02265852d0aacad5c97d70f078e3a}}\\
\multicolumn{2}{@{}l@{}}{\quad internal \texttt{c40d27be152217bd0d862bf54c4a957ab7bcad22e6b2f7747cf7462e9bc7a397}}\\
\multicolumn{2}{@{}l@{}}{\quad module \texttt{99840c239eb60c3db3a1093827f3c9849c5d010debec70eb13661a87e0418fe0}}\\
\addlinespace
\path{liu9-psi-reduction.json} (\path{liu9_psi_reduction_audit.py}) & independent re-derivation of the chain; \texttt{PROVED}\\
\multicolumn{2}{@{}l@{}}{\quad file \texttt{9fb5766f1c3110cce9b231d2f2f1b7c1eec67bc3d6ac7602be50fc72b00c07a3}}\\
\multicolumn{2}{@{}l@{}}{\quad internal \texttt{405eea478f2146fde762f749dd4a8f6986be79f0be47f56cb2cc24daf64b2af0}}\\
\multicolumn{2}{@{}l@{}}{\quad module \texttt{18602dc7c8599ab6752583535ed5809df571deb8d8851e4accb4d07362e927bb}}\\
\addlinespace
\path{liu9-h2-twovar.json} (\path{liu9_h2_twovar_lemmas.py}) & $R\ge0$, $A\ge0$, Theorem~\ref{thm:twovar}; \texttt{CERTIFIED\_PROVED}\\
\multicolumn{2}{@{}l@{}}{\quad file \texttt{a028ff17ae8492211a2ef9de5721de7f3a0afd8bd45bf7d8d9b77b7c90d2e6eb}}\\
\multicolumn{2}{@{}l@{}}{\quad internal \texttt{4ddbf0f38ca881dde204f086076424617edcdbe6a2826dfa945ef963ebdf70c6}}\\
\multicolumn{2}{@{}l@{}}{\quad module \texttt{2afe8e242aa5f506c6d54967a133fe246a7e5cb8bad4961b9860e93b86f67163}}\\
\addlinespace
\path{liu9-h2-boundary.json} (\path{liu9_h2_boundary_layer.py}) & $\Phi\ge0$, Theorem~\ref{thm:phi}; \texttt{PROVED}\\
\multicolumn{2}{@{}l@{}}{\quad file \texttt{e4d6d96df3ae435800cef5739fb1ae45b98febf539ab418df116ceb03bd481c9}}\\
\multicolumn{2}{@{}l@{}}{\quad internal \texttt{389804b249cd8b16a47f3c6ab64a70ad4a9afe78097b4a8c6739759804c9eaaa}}\\
\multicolumn{2}{@{}l@{}}{\quad module \texttt{fcb3ed55e8bd9ad6d2f180f5ffa79f192dacc940a5c7506eca2a85aa612905e6}}\\
\addlinespace
\path{liu9-h2-phi-audit.json} (\path{liu9_h2_phi_audit.py}) & independent audit of $\Phi\ge0$; \texttt{AUDIT-PASSED}\\
\multicolumn{2}{@{}l@{}}{\quad file \texttt{f2359fa3b3b08bab29965253d777dcf32e4b1aee1e0bacbfb0924b86ac163996}}\\
\multicolumn{2}{@{}l@{}}{\quad internal \texttt{a515102c1807849b606bf4a681bd0a67173445b1a22610d397bb779440222cab}}\\
\multicolumn{2}{@{}l@{}}{\quad module \texttt{b1d9ed3d9da9e7c99c026843b683447809fb38d28dba4c49a087729bf66cb50c}}\\
\addlinespace
\path{liu9-h2-general-lift.json} (\path{liu9_h2_general_lift.py}) & Theorem~\ref{thm:lift}; \texttt{PROVED}\\
\multicolumn{2}{@{}l@{}}{\quad file \texttt{eb4874d39904593bddc3f2ae638d1bad6756a19d4ebf6e06a60427b6bb7ac1ce}}\\
\multicolumn{2}{@{}l@{}}{\quad internal \texttt{3ff1e797cd5dbcc4273d9f002b9e540b58829d237b2d344f266e0bea46c02bc8}}\\
\multicolumn{2}{@{}l@{}}{\quad module \texttt{b265e726d110098a09c00de331c4433454973ba741966b3018b77afcc225b5c3}}\\
\addlinespace
\path{liu9-h2-mixture-theorem.json} (\path{liu9_h2_mixture_theorem.py}) & Theorem~\ref{thm:mixture}; \texttt{PROVED}\\
\multicolumn{2}{@{}l@{}}{\quad file \texttt{329f7e2d71af8cd78d1a921c72b9ae4d05b71113134eb3341d69932f19ea24b3}}\\
\multicolumn{2}{@{}l@{}}{\quad internal \texttt{b108b781224dff601ffcf3f562703d3b620e7da6d89cf6d000b6ccbd06a4773b}}\\
\multicolumn{2}{@{}l@{}}{\quad module \texttt{3098a1ca30a0f16582df02e1fb7dfd8fa27970c031167125e696e58ac43dfbab}}\\
\addlinespace
\bottomrule
\end{tabular}
\caption{The certificate chain.  For each artifact: SHA-256 of the file, the
internal digest, and the module source.  All $21$ values are pinned in
\texttt{PROGRESS.md} and asserted by the replay driver.}\label{tab:artifacts}
\end{table}

\paragraph{Replay.}
From \path{math/},
\begin{center}
\texttt{nice -n 19 ./.venv/bin/python -I -B uc/verification/replay\_h2\_chain.py}
\end{center}
re-runs each of the seven modules twice from the repository root under the
same interpreter, requires every run to reproduce the on-disk artifact byte
for byte, recomputes the three hashes of every artifact from disk in the
module's own canonicalisation, checks the cross-artifact pins, and asserts
every hash against the values pinned in the repository ledger
\path{PROGRESS.md} (which contains the same $21$ values as
Table~\ref{tab:artifacts}); the exit status is nonzero on any mismatch.
The modules use python-flint (Arb ball arithmetic \cite{Joh17}; $400$ bits
for the certificates of Sections~\ref{sec:twovar}--\ref{sec:mixture} and the
audit, $480$ for the reduction, $320$ for its independent re-derivation) and
sympy, single-threaded; the constants are
never re-rounded through floating point, and the sole external inputs are
the two audited transcriptions \path{uc/liu9_objective.py} (Liu's objective)
and \path{uc/liu9_binding.py} (his equations (87)--(90)), both pinned by the
lift and mixture artifacts.  Each module carries mutations --- deliberate
errors in the statement or the certificate --- and refuses to certify when
any of them is active; all $42$ mutations across the chain fire.

\paragraph{What is proved and what is not.}
Theorems~\ref{thm:reduction}, \ref{thm:twovar}, \ref{thm:phi},
\ref{thm:lift}, \ref{thm:mixture} are \lbl{proved} with
\lbl{machine-verified} finite components (identities, enclosures, covers);
Theorem~\ref{thm:uc} is \lbl{proved} by the elementary argument of
Section~\ref{sec:prop3}; Theorem~\ref{thm:value} is
\lbl{machine-verified}.  Two blank-context adversarial reviews (one of the
$\Phi$ certificate, one of the lift and mixture artifacts and of the
argument of Section~\ref{sec:prop3}) found no soundness defect in the final
versions; the first version of the lift artifact was refuted by the
reviewer at the truncated constants (Remark~\ref{rem:exact}) and repaired.
The result has not been refereed outside the repository; that is recorded
as the adopted next step.  Nothing here concerns the conjectured value
$\tfrac12$: Lemma~\ref{lem:sharp} shows that $c'$ is the best constant this
particular pair of protocols at $\beta^\ast$ can give, and improving it
requires a different protocol family.

\paragraph{Why interval subdivision alone could not have closed
Hypothesis~2.}
$R$ has a nondegenerate interior zero (Lemma~\ref{lem:tight}); a function
that touches zero cannot be proved nonnegative by bisection at any depth.
The stratification of Section~\ref{sec:twovar} handles the zero by exact
algebra plus a Taylor bound, and the identity of Section~\ref{sec:phi}
converts the two-variable $B\ge0$, which vanishes on the entire diagonal
(and on both axes), into the one-variable condition $\kappa\ge2$ with a
uniform margin.  Both are
what makes the dimension-free reduction of Section~\ref{sec:reduction}
usable: Hypothesis~2, a nine-parameter statement, and its general-measure
form are decided by two inequalities on the unit square.

\begin{thebibliography}{99}
\bibitem{AHS22} R.~Alweiss, B.~Huang, M.~Sellke, Improved lower bound for the
union-closed sets conjecture, arXiv:2211.11731 (2022).
\bibitem{Cam22} S.~Cambie, Better bounds for the union-closed sets conjecture
using the entropy approach, arXiv:2212.12500 (2022).
\bibitem{CL22} Z.~Chase, S.~Lovett, Approximate union closed conjecture,
arXiv:2211.11689 (2022).
\bibitem{Gil22} J.~Gilmer, A constant lower bound for the union-closed sets
conjecture, arXiv:2211.09055 (2022).
\bibitem{Joh17} F.~Johansson, Arb: efficient arbitrary-precision
midpoint-radius interval arithmetic, IEEE Trans.\ Computers 66 (2017),
1281--1292.
\bibitem{Liu23} J.~Liu, Improving the lower bound for the union-closed sets
conjecture via conditionally IID coupling, arXiv:2306.08824v1 (2023).
\bibitem{Saw22} W.~Sawin, An improved lower bound for the union-closed set
conjecture, arXiv:2211.11504 (2022).
\bibitem{Yu23} L.~Yu, Dimension-free bounds for the union-closed sets
conjecture, Entropy 25 (2023), 767.
\end{thebibliography}

\end{document}
